Este trabalho propõe o desenvolvimento e validação da Fase 1 de um protótipo modular de Realidade Virtual Imersiva para treinamento de segurança em canteiros de obras, com foco na identificação e mitigação de riscos relacionados a trabalhos em altura e uso de Equipamentos de Proteção Individual (EPIs), no contexto brasileiro. O sistema será desenvolvido em Unity 3D com Meta XR SDK para dispositivos como o Oculus Quest 3. A metodologia envolve um estudo piloto com 30 participantes divididos em grupo experimental (Realidade Virtual) e grupo controle (treinamento tradicional por vídeo instrucional). A avaliação utilizará instrumentos padronizados: testes de conhecimento pré e pós-treinamento, questionários de motivação e engajamento, \textit{System Usability Scale} (SUS) e \textit{iGroup Presence Questionnaire} (IPQ). Este estudo contribui com o estabelecimento de um protocolo de avaliação adaptado ao contexto brasileiro, sob as diretrizes das Normas Regulamentadoras nacionais, fornecendo os primeiros indicadores nacionais para esse tipo de intervenção.

% Separe as palavras-chave por ponto
\palavraschave{Realidade Virtual. Segurança do Trabalho. Construção Civil. Treinamento Imersivo. NR-18.}
